% Exact dyadic finite-ramification groups and their tower consequence.
\subsection{At coefficient two the fixed-ramification obstruction is real}

\begin{result}{E.155}{Proved}{All dyadic weights and the exact real obstruction}
Put \(S_2=\{2,\infty\}\). For every integer \(w\) and every \(n\geq1\),
restriction is an isomorphism
\[
 H^2(G_{\mathbb Q,S_2},\mathbb Z/2^n(w))
       \xrightarrow{\ \sim\ }
 H^2(G_{\mathbb R},\mathbb Z/2^n(w)).
 \tag{E.155.1}
\]
Both finite groups have order two. With continuous integral coefficients,
\[
 H^2_{\rm cts}(G_{\mathbb Q,S_2},\mathbb Z_2(w))
 \xrightarrow{\ \sim\ }H^2(G_{\mathbb R},\mathbb Z_2(w))
 =
 \begin{cases}
   \mathbb Z/2,&w\text{ even},\\
   0,&w\text{ odd}.
 \end{cases}
 \tag{E.155.2}
\]
For \(w\geq2\) and any finite \(S\supset S_2\), real restriction
and the actual tame residues give an isomorphism
\[
 H^2_{\rm cts}(G_{\mathbb Q,S},\mathbb Z_2(w))
 \xrightarrow{\ \sim\ }
 H^2(G_{\mathbb R},\mathbb Z_2(w))
 \oplus\!
 \bigoplus_{q\in S\setminus S_2}
      \mathbb Z_2/(q^{w-1}-1).
 \tag{E.155.3}
\]
In particular A.22 gives the canonical absolute description
\[
 \begin{split}
 H^2_{\rm cts}(G_{\mathbb Q},\mathbb Z_2(w))
 &\xrightarrow{\ \sim\ }
 H^2(G_{\mathbb R},\mathbb Z_2(w))\oplus R_2/d_wR_2,\\
 R_2&=c_0(\{q:q\ne2\},\mathbb Z_2),\qquad
 d_w(y)_q=(q^{w-1}-1)y_q .
 \end{split}
 \tag{E.155.4}
\]
The denominator is the actual image, so (E.155.4) retains
the possible non-Hausdorff and derived-limit contribution.

Consequently a pro-two class-\(r\) section satisfying the
ordinary boundary matching hypothesis of B.25 has a global
integral class-\((r+1)\) lift if and only if its real restriction
has such a lift. If it is geometrically conjugate at the real
place to a specified ordinary point section, a global lift exists
and may be chosen to retain that real point class. Simultaneous
matching at the finite support primes remains the independent
first-cohomology condition; this theorem does not assert its
consistency.
\end{result}
\begin{proof}
Write \(G_2=G_{\mathbb Q,S_2}\), \(T=\mathbb Z_2(w)\), and
\(T_n=T/2^nT\). Every \(T_1\) is the trivial module
\(\mathbb F_2=\mu_2\), independently of \(w\).
Purity, the comparison with \(G_2\), and Kummer theory give
\[
 \mathbb Q^\times/\mathbb Q^{\times2}
 \longrightarrow\bigoplus_{q\ne2}\mathbb F_2
 \longrightarrow H^2(G_2,\mu_2)
 \longrightarrow\operatorname{Br}(\mathbb Q)[2].
 \tag{E.155.5}
\]
The first map is the vector of valuations modulo two, and is
surjective. These are the comparison and supports maps of
\citet[II.2.9, pp.~170--171; II.1.7(b), pp.~151--152]{CohMilneADT}.
Thus the last map in (E.155.5) is injective. Its image is
unramified at every odd prime. Brauer reciprocity therefore
identifies a class in its image by its real invariant: its
two-adic invariant is the negative of that invariant.
The class \(\kappa(-1)\cup\kappa(-1)\) belongs to
\(H^2(G_2,\mu_2)\) and has nonzero real restriction.
Hence restriction to the real group is an isomorphism of
groups of order two. The Brauer injection and invariant sum
used here are \citet[II.2.1--2.2, pp.~163--165]{CohMilneADT}.

We next verify the degree-one endpoint needed in a coefficient
induction. For every \(w,n\),
\[
             H^1(G_2,T_n)\longrightarrow H^1(G_{\mathbb R},T_n)
             \quad\hbox{is surjective}.
 \tag{E.155.6}
\]
Let \(c\) denote complex conjugation. If \(w\) is odd,
\(c\) acts as minus one and the target is \(T_n/2T_n\).
The integral cocycle
\[
                  g\longmapsto\frac{\chi_2(g)^w-1}{2}
 \tag{E.155.7}
\]
is defined for every integer \(w\), is unramified outside two,
and has value \(-1\) at \(c\). Its reduction generates the
target. If \(w\) is even, \(c\) acts trivially and the target
is \(T_n[2]\). The equivariant inclusion
\(\mathbb F_2\hookrightarrow T_n\), \(1\mapsto2^{n-1}\),
sends the quadratic character \(\kappa(-1)\) to a cocycle
whose value at \(c\) is \(2^{n-1}\). This generates the
target in the even case. These are ordinary degree-one
cohomology groups; the real Tate convention has the same
groups in positive degree.

Use the literal coefficient sequence
\[
 0\longrightarrow T_1
 \xrightarrow{\,2^n\,}T_{n+1}
 \longrightarrow T_n\longrightarrow0.
 \tag{E.155.8}
\]
All maps use one fixed Tate module \(T\), so they respect
the full two-adic cyclotomic action. Compare the degree-one
through degree-three exact sequences for \(G_2\) and
\(G_{\mathbb R}\). The degree-two vertical maps for \(T_1\)
and \(T_n\) are isomorphisms by the preceding calculation
and induction. The degree-three vertical map for \(T_1\)
is an isomorphism by
\citet[I.4.10(c), p.~57]{CohMilneADT}.
Finally (E.155.6) makes the degree-one vertical map for
\(T_n\) surjective.

For completeness these facts imply the middle isomorphism
without suppressing a Bockstein. If a global class for
\(T_{n+1}\) has zero real image, its projection for \(T_n\)
is zero. Write it as the image of a class for \(T_1\).
The real image of that latter class is the boundary of a
real degree-one class for \(T_n\); lift the degree-one class
by (E.155.6) and use the degree-two isomorphism for \(T_1\).
This proves injectivity. Conversely, project a real degree-two
class for \(T_{n+1}\) to \(T_n\) and lift it globally. Its
boundary for \(T_1\) is zero because its real image is zero
and the degree-three map is injective. It therefore lifts
globally to \(T_{n+1}\). The remaining difference comes
from \(H^2(G_{\mathbb R},T_1)\) and can be corrected globally.
This proves surjectivity and completes (E.155.1).

The periodic resolution for the group \(\langle c\rangle\)
gives
\[
 H^2(G_{\mathbb R},T_n)=
 \begin{cases}
 T_n/2T_n,&w\text{ even},\\
 T_n[2],&w\text{ odd}.
 \end{cases}
\]
Reduction from exponent \(n+1\) to \(n\) is the identity
on the displayed order-two group in even weight, and is
zero in odd weight. This distinction is essential:
the finite odd-weight groups are nonzero although their
integral inverse limit is zero.

For fixed finite ramification, the first and second finite
cohomology groups are finite
\citep[I.4.9--4.10, pp.~56--57]{CohMilneADT}.
Their inverse systems are therefore Mittag--Leffler.
Jannsen's continuous-cohomology exact sequence
\citep[(2.1) and Theorem~2.2, p.~215]{CohJannsen1988}
identifies the integral second group with the inverse
limit of the finite second groups. It also shows that
for any finite \(S\supset S_2\), the integral third
restriction
\[
 H^3_{\rm cts}(G_{\mathbb Q,S},T)
       \longrightarrow H^3(G_{\mathbb R},T)
 \tag{E.155.9}
\]
is an isomorphism: apply the same argument to the finite
degree-three isomorphisms of Milne I.4.10(c).
This proves (E.155.2) with all inverse-limit terms accounted for.

Now suppose \(w\geq2\). Integral purity for adding finitely
many odd primes to \(S_2\) has exact segment
\[
 \begin{split}
 0\longrightarrow H^2_{\rm cts}(G_2,T)
 &\longrightarrow H^2_{\rm cts}(G_{\mathbb Q,S},T)
 \longrightarrow
   \bigoplus_{q\in S\setminus S_2}H^1(\mathbb F_q,\mathbb Z_2(w-1))\\
 &\longrightarrow H^3_{\rm cts}(G_2,T)
 \longrightarrow H^3_{\rm cts}(G_{\mathbb Q,S},T).
 \end{split}
 \tag{E.155.10}
\]
This is the finite-set form of the coefficient-compatible
purity triangle proved in A.22. The zero on the left follows
from \(H^0(\mathbb F_q,\mathbb Z_2(w-1))=0\).
The last arrow is an isomorphism by (E.155.9), since
both real restrictions are the identity on the same group.
Also
\[
 H^1(\mathbb F_q,\mathbb Z_2(w-1))
                  =\mathbb Z_2/(q^{w-1}-1).
 \]
Thus (E.155.10) is a short exact sequence at its first three
nonzero terms. Its kernel is identified isomorphically with
the real group by (E.155.2). Real restriction together with
the residues is consequently the isomorphism (E.155.3):
given residues, lift them and then correct its real component
by the unique kernel class. The same argument applied to
the strict absolute sequence A.22 gives (E.155.4).

Finally B.25 places the next obstruction of a boundary-matched
section in the image of
\(H^2_{\rm cts}(G_2,\mathbb Z_2(r+1))^{b_{r+1}}\).
The restriction of this class at infinity is precisely
the real lifting obstruction. Equation (E.155.2) proves
the stated equivalence. If the preceding real class is a
point class, its obstruction vanishes because that point
has a full section. Choose a global lift and align its
lower real restriction with that of the specified point.
The difference of these real lifts is a degree-one class
in \(\mathbb Z_2(r+1)^{b_{r+1}}\). Integral restriction in
degree one is surjective: its real target is zero in even
weight, and in odd weight is generated by (E.155.7).
A global twist therefore retains the real point class.
This last correction concerns the real condition alone;
retaining all finite boundary conditions is still a joint
localization-image problem, with all required aligning
frames retained.
\end{proof}

\begin{result}{E.156}{Proved}{The exact dyadic first-cohomology parameters and localization}
For \(w\geq2\), define
\[
 a_w=
 \begin{cases}
  1,&w\text{ odd},\\
  2+v_2(w),&w\text{ even}.
 \end{cases}
 \qquad
 \eta_w(g)=\frac{\chi_2(g)^w-1}{2^{a_w}}.
 \tag{E.156.1}
\]
Inflation identifies global continuous first cohomology with
that over \(G_{\mathbb Q,S_2}\), and
\[
 H^1_{\rm cts}(G_{\mathbb Q},\mathbb Z_2(w))\simeq
 \begin{cases}
  \mathbb Z_2\oplus\mathbb Z/2,&w\text{ odd},\\
  \mathbb Z/2^{\,2+v_2(w)},&w\text{ even}.
 \end{cases}
 \tag{E.156.2}
\]
The torsion is generated by \(\eta_w\). In odd weight
its real restriction is an isomorphism onto
\(H^1(G_{\mathbb R},\mathbb Z_2(w))=\mathbb Z/2\).
The kernel of real restriction is a free rank-one
\(\mathbb Z_2\)-module. In even weight the real target is zero.

At the coefficient prime,
\[
 H^1_{\rm cts}(G_{\mathbb Q_2},\mathbb Z_2(w))
            \simeq\mathbb Z_2\oplus\mathbb Z/2^{a_w}.
 \tag{E.156.3}
\]
For odd \(w\geq3\), localization from the global group
to (E.156.3) is an isomorphism. For even \(w\) its
image is exactly the local torsion subgroup.
In particular an odd-weight correction prescribed at two
determines its real correction as well; the real and
two-adic conditions are not independent parameters.
\end{result}
\begin{proof}
For \(q\ne2\), the inertia of an integral weight-\(w\)
one-cocycle is annihilated by \(q^{w-1}-1\). The coefficient
module is torsion free and this integer is nonzero.
Thus inertia vanishes. The cocycle, as well as every
coboundary between such cocycles, descends to \(G_2\).
This is also the finite-ramification calculation of A.20.

The full cyclotomic image is
\(\mathbb Z_2^\times=\{1,-1\}\times(1+4\mathbb Z_2)\).
For odd \(w\), the element \(-1\) gives valuation one for
\(\chi^w-1\). For even \(w\), the generator \(5\) of
\(1+4\mathbb Z_2\) gives
\[
                 v_2(5^w-1)=2+v_2(w).
 \tag{E.156.4}
\]
For instance this identity follows from the odd-exponent
geometric sum and successive doubling; after the first
step every factor \(5^{2^ju}+1\) has valuation one.
It follows that
\[
 H^0(G_2,\mathbb Q_2/\mathbb Z_2(w))
        =\frac{2^{-a_w}\mathbb Z_2}{\mathbb Z_2}.
 \tag{E.156.5}
\]
There are no invariants in \(\mathbb Q_2(w)\).
The coefficient boundary therefore identifies (E.156.5)
with the entire torsion in \(H^1(G_2,\mathbb Z_2(w))\),
and gives exactly the cocycle \(\eta_w\).

Set \(M_n=\mathbb Z/2^n(w)\).
The finite global Euler characteristic
\citep[I.5.1, p.~67; I.5.2(c), p.~68]{CohMilneADT},
together with E.155, gives
\[
 \begin{split}
 |H^0(G_2,M_n)|&=2^{\min(n,a_w)},\\
 |H^1(G_2,M_n)|
  &=|H^0(G_2,M_n)|\,|H^2(G_2,M_n)|
           \frac{|M_n|}{|H^0(G_{\mathbb R},M_n)|}\\
  &=
  \begin{cases}
   2^{n+1},&w\text{ odd},\\
   2^{\min(n,a_w)+1},&w\text{ even}.
  \end{cases}
 \end{split}
 \tag{E.156.6}
\]
Here \(H^0(G_{\mathbb R},M_n)\) is ordinary invariants,
of size two in odd weight and \(2^n\) in even weight,
as required in the Euler formula; it is not Tate \(H^0\).

The integral first group is compact, being the inverse
limit of its finite first-cohomology groups; the
\(\varprojlim^1H^0\) term vanishes by finiteness.
Its quotient by two is finite, by the coefficient exact
sequence, so compact Nakayama makes it a finitely
generated \(\mathbb Z_2\)-module. The same sequence gives
\[
 0\longrightarrow H^1(G_2,\mathbb Z_2(w))/2^n
 \longrightarrow H^1(G_2,M_n)
 \longrightarrow H^2_{\rm cts}(G_2,\mathbb Z_2(w))[2^n]
 \longrightarrow0.
 \tag{E.156.7}
\]
For odd weight the last group is zero, and for even
weight it has order two. Equations (E.156.5)--(E.156.7)
prove (E.156.2): the free rank is one in odd weight and
zero in even weight. Restriction of \(\eta_w\) at complex
conjugation is \(-1\) in odd weight. Since the real first
group then has order two, this proves the asserted
real splitting and its torsion-free rank-one kernel.

The same cyclotomic image occurs over \(\mathbb Q_2\).
Thus (E.156.5) also gives the local torsion, with the
same generator. Local Tate duality makes
\[
 H^2_{\rm cts}(G_{\mathbb Q_2},\mathbb Z_2(w))
 \quad\hbox{dual to}\quad
 H^0(G_{\mathbb Q_2},\mathbb Q_2/\mathbb Z_2(1-w)),
 \tag{E.156.8}
\]
a finite group for \(w\geq2\).
Passing from finite coefficients to (E.156.8) is valid
because local finite first cohomology is finite.
The local Euler characteristic
\citep[I.2.3, p.~28; I.2.8, p.~31]{CohMilneADT}
then gives free rank one in degree one, proving
(E.156.3). Equivalently, in the finite Euler formula the
factor \(2^n\) supplies that rank, while the two finite
invariant groups account for torsion and the finite
second group. This also proves finite generation as above.

It remains to check the integral index of localization
in odd weight; equal ranks alone would not suffice.
By E.155 and (E.156.7), reduction modulo two identifies
the global first group modulo two with
\[
 H^1(G_2,\mathbb F_2)
       =\langle\kappa(-1),\kappa(2)\rangle.
 \tag{E.156.9}
\]
This is Kummer theory for the units of \(\mathbb Z[1/2]\).
The torsion generator reduces to \(\kappa(-1)\):
for odd \(w\), \((\chi^w-1)/2\) modulo two is the
quadratic character of \(\mathbb Q(i)\).
Choose a global integral class reducing to \(\kappa(2)\).
It has zero real restriction, because reduction on the
real degree-one group is an isomorphism in odd weight.
It therefore generates the free real kernel modulo two.

The local Kummer classes of \(-1\) and \(2\) are linearly
independent in \(\mathbb Q_2^\times/\mathbb Q_2^{\times2}\):
the latter has odd valuation, while the former is a
nonsquare unit. Reduction injects the local integral
first group modulo two into this finite Kummer group.
Consequently the image of the chosen global free generator
is not in the span of local torsion modulo two. Its
coefficient in the local free rank-one quotient is a unit.
Localization is already an isomorphism on torsion, so
it is an isomorphism on the entire odd-weight groups.
For even weight the global group consists only of the
same torsion (E.156.5), proving the final assertion.
\end{proof}
